% =====================================================================
%  Constraint Geometry and Domain Physics
%  A conceptual synthesis note (methodological, non-theorem-bearing)
%
% =====================================================================
\documentclass[11pt,a4paper]{article}

\usepackage[T1]{fontenc}
\usepackage[utf8]{inputenc}
\IfFileExists{lmodern.sty}{\usepackage{lmodern}}{}
\usepackage{amsmath,amssymb,amsthm}
\usepackage{mathtools}
\usepackage{booktabs}
\usepackage[margin=1in]{geometry}
\IfFileExists{microtype.sty}{\usepackage[protrusion=true,expansion=false]{microtype}}{}
\usepackage{enumitem}
\usepackage[hidelinks]{hyperref}

% --- environments -----------------------------------------------------
\theoremstyle{definition}
\newtheorem{definition}{Definition}[section]
\theoremstyle{plain}
\newtheorem{proposition}[definition]{Proposition}
\newtheorem{lemma}[definition]{Lemma}
\theoremstyle{remark}
\newtheorem{remark}[definition]{Remark}
\newtheorem{observation}[definition]{Observation}
\newtheorem{interpretation}[definition]{Structural Interpretation}
\newtheorem{safeguard}[definition]{Safeguard}

% --- macros -----------------------------------------------------------
\newcommand{\Ccrit}{C_{\mathrm{crit}}}
\newcommand{\Rresp}{R}
\newcommand{\lin}{{(1)}}
\newcommand{\nonlin}{{(\ge 2)}}

\title{\bfseries Constraint Geometry and Domain Physics:\\
A Two-Layer Architecture for Binary Partition Diagnostics}

\author{Elias De Jes\'us\\
\small Independent Researcher\\
\small ORCID: 0009-0007-0190-9143\\
\small \texttt{dejesuselias10@gmail.com}}

\date{July 2026 --- Conceptual Synthesis Note\\[2pt]
\small DOI: \href{https://doi.org/10.5281/zenodo.21695430}{10.5281/zenodo.21695430}}

\begin{document}
\maketitle

\begin{abstract}
\noindent
This note clarifies the mature architecture of a binary-partition research
programme without asserting a new theorem or a new physical law. The organizing
claim is a division of labor: constraint geometry together with domain physics
yields a physical diagnostic. The constraint layer is the conserved binary
partition $a+b=1$ with Thales altitude $h=\sqrt{ab}$, asymmetry $\chi=a-b$ and
response $R=1/(1-\chi^{2})=1/(4h^{2})$, together with intrinsic landmarks such
as the crossing $h=a/b$, equivalent on the partition to $a=b^{3}$. The domain layer supplies the identity of
the two sectors, the reason they form a normalized partition, the map
$a/b=F(X)$ between a physical variable and the partition coordinates, the
trajectory through the manifold, and the empirical meaning of any landmark. The
perturbative decomposition of the Einstein tensor into linear and higher-order
parts is developed as the primary example: general relativity motivates the
schematic sector scaling $a/b\sim C$ with $C$ the compactness, and the
intrinsic crossing then reads $C(1+C)^{2}=1$. The note emphasizes that the
geometric landmark is fixed in partition space while its physical coordinate
$X_{*}$ moves with the embedding $F$, so that an unjustified choice of $F$ can
manufacture any desired threshold. A five-step protocol and an explicit
epistemic tiering are given, and an illustration in flat $\Lambda$CDM
separates two failure modes: one landmark there is derived from the dynamics
but unoccupied, another occupied but provably mechanism-free.
\end{abstract}

\noindent\textbf{Keywords:} constraint geometry; binary partition; Thales
altitude; domain embedding; intrinsic relational landmark; compactness
diagnostic; perturbative general relativity; epistemic tiering.

\tableofcontents

% =====================================================================
\section{Introduction}
% =====================================================================

\subsection{Purpose}

A conserved two-sector decomposition appears in many parts of mathematics and
physics: a probability allocation, a mass fraction, an energy budget, a density
partition, a ratio of operator norms. Whenever two positive quantities are
normalized so that they sum to one, the same small piece of geometry becomes
available: a bounded state coordinate, a geometric mean, an asymmetry measure,
a response measure, and a short list of relational landmarks fixed by the
constraint alone.

That availability is not a physical result, and the central purpose of this
note is to say precisely what it is instead. The programme summarized here has
at various points been described in language that invited a stronger reading
than intended --- as though partition geometry were generating domain physics.
The mature formulation is a division of labor,
\begin{equation}
\boxed{\;
\text{constraint geometry}
\;+\;
\text{domain physics}
\;\longrightarrow\;
\text{physical diagnostic}\;}
\label{eq:architecture}
\end{equation}
in which the two layers contribute different and non-interchangeable things.
The note is a synthesis and clarification of results developed in earlier work
\cite{TwoFace,EffPot,Newtonian2Body,IntrinsicCubic,CubicThreshold,ReverseThales};
it introduces no new theorem and no new physical law.

The central interpretive statement is the following.

\begin{quote}
\emph{Constraint geometry determines what relational structures and landmarks
are available; domain physics determines which of them are physically realized
and what they mean.}
\end{quote}

\subsection{What the two layers supply}

The constraint layer supplies the invariant relational structure
\begin{equation}
a+b=1,\qquad
h=\sqrt{ab},\qquad
\chi=a-b,\qquad
\Rresp=\frac{1}{1-\chi^{2}} ,
\label{eq:coords}
\end{equation}
together with the coordinate identities among these quantities and any
relational condition expressible in them alone.

Domain physics supplies, separately and independently:
\begin{enumerate}[label=(\roman*),leftmargin=*]
\item the physical meaning of the two sectors;
\item the reason those sectors form a genuine normalized partition rather than
      an imposed one;
\item the map between a physical variable and the partition coordinates;
\item the trajectory through the partition manifold, if the system evolves;
\item the empirical meaning, if any, of a geometric landmark.
\end{enumerate}

Neither layer produces a diagnostic alone. Constraint geometry without an
embedding is a list of identities; an embedding without independent physical
justification is a change of variables chosen to produce a desired answer.

\subsection{Non-claims}
\label{sec:nonclaims}

This note does \emph{not} claim:
\begin{enumerate}[label=(N\arabic*),leftmargin=*]
\item that constraint geometry derives the Einstein field equations, or any
      other domain field equation;
\item that the intrinsic cubic is a universal law of gravity, a fundamental
      constant of nature, or a universal physical phase transition;
\item that nonlinear gravity begins suddenly at the cubic point, or that the
      point marks a sharp physical onset of any kind;
\item that two domains share dynamics because they share a partition;
\item that a partition-space landmark has physical meaning prior to an
      independently justified embedding;
\item that the numerical value of any physical threshold obtained here is
      fixed independently of the schematic normalization of the sector ratio;
\item that any of the illustrations in \S\ref{sec:crossdomain} constitutes
      empirical validation;
\item that proximity of an observed value to a landmark is evidence, when the
      neighbourhood of that landmark contains generic constants within
      measurement error.
\end{enumerate}

% =====================================================================
\section{Binary Partition Constraint Geometry}
\label{sec:geometry}
% =====================================================================

\begin{definition}[Partition data]
Let $a+b=1$ with $a,b>0$. Define the Thales altitude $h=\sqrt{ab}$, the
asymmetry $\chi=a-b\in(-1,1)$, the response $\Rresp=1/(1-\chi^{2})$, the
rapidity $\xi=\operatorname{arctanh}\chi$, and the partition ratio
$r=a/b\in(0,\infty)$.
\end{definition}

Since $a=\tfrac12(1+\chi)$ and $b=\tfrac12(1-\chi)$, one has
$ab=(1-\chi^{2})/4$, hence the exact constraint identity
\begin{equation}
4h^{2}+\chi^{2}=1,
\qquad
\Rresp=\frac{1}{1-\chi^{2}}=\frac{1}{4h^{2}} .
\label{eq:thales}
\end{equation}
Equation~\eqref{eq:thales} is the analytic content of the Thales semicircle:
the altitude and the half-asymmetry are the legs of a right triangle on a unit
diameter. It is exact, parameter-free, and carries no physics.

\begin{observation}[Bounded state, unbounded response]
The state coordinate is bounded while the response and rapidity are not:
$\chi\in(-1,1)$, $\Rresp\in[1,\infty)$, $\xi\in\mathbb{R}$. A divergence in
$\Rresp$ therefore records approach to a partition boundary, in whatever
coordinate the domain provides, and is not by itself evidence of a physical
singularity \cite{TwoFace}.
\end{observation}

\subsection{Intrinsic relational landmarks}

A relational landmark is a locus on the partition manifold specified by a
condition written in the coordinates \eqref{eq:coords} alone. The landmark used
throughout this note is the altitude--ratio crossing.

\begin{proposition}[Altitude--ratio crossing; exact]
\label{prop:crossing}
On $a+b=1$ with $a,b>0$,
\begin{equation}
h=\frac{a}{b}
\quad\Longleftrightarrow\quad
a=b^{3}
\quad\Longleftrightarrow\quad
b^{3}+b-1=0 .
\label{eq:crossing}
\end{equation}
The cubic $b^{3}+b-1$ is strictly increasing on $[0,1]$ with a sign change, so
the crossing occurs at exactly one interior state, and
\begin{equation}
b_{*}=0.682327803828019\ldots,\qquad
a_{*}=b_{*}^{3}=0.317672196171981\ldots,
\end{equation}
\begin{equation}
h_{*}=b_{*}^{2}=r_{*}=\frac{a_{*}}{b_{*}}
=0.465571231876768\ldots=\psi-1,
\end{equation}
where $\psi=1/b_{*}=1.465571\ldots$ satisfies $\psi^{3}=\psi^{2}+1$. In
supergolden coordinates $(a_{*},h_{*},b_{*})=(\psi^{-3},\psi^{-2},\psi^{-1})$,
with $\chi_{*}=a_{*}-b_{*}=-0.3646556\ldots$ and
$\Rresp_{*}=\psi^{4}/4=1.1533676\ldots$
\end{proposition}

\begin{proof}
If $h=a/b$ then $ab=a^{2}/b^{2}$, so $ab^{3}=a^{2}$ and, dividing by $a>0$,
$b^{3}=a$. Conversely $a=b^{3}$ gives $h=\sqrt{b^{3}\cdot b}=b^{2}$ and
$a/b=b^{2}$. Substituting $a=1-b$ gives $b^{3}+b-1=0$; the derivative
$3b^{2}+1>0$ gives uniqueness. The identity $\psi-1=\psi^{-2}$ follows from
$\psi^{3}=\psi^{2}+1$ on division by $\psi^{2}$; the closed forms for
$h_{*}$ and $\Rresp_{*}$ then follow from \eqref{eq:thales}. These facts are
established in \cite{IntrinsicCubic}.
\end{proof}

\begin{lemma}[Sign law; exact]
\label{lem:sign}
On the partition, $h>a/b\iff b^{3}>a\iff b>b_{*}$, and $h<a/b\iff b<b_{*}$.
The crossing is therefore not a tangency: the two functions cross
transversally, in the sense that the difference changes sign there.
\end{lemma}

\begin{remark}[What the landmark is, and is not]
\label{rem:landmarkstatus}
Proposition~\ref{prop:crossing} is a statement about two functions of one
constrained variable. It says that the geometric mean of the two sectors and
their ratio agree at exactly one interior state. It does not say that anything
happens there. Nothing in \eqref{eq:crossing} refers to a physical system, and
the number $0.465571\ldots$ is at this stage an algebraic constant of the
partition, on the same footing as the apex value $h=\tfrac12$ at $\chi=0$.
\end{remark}

\begin{remark}[A caution on levels]
\label{rem:levels}
The altitude is a geometric mean of the two fractions --- an amplitude-level
object --- while the ratio $a/b$ compares them directly. Setting the two equal
is a partition-intrinsic relational condition; it is not a balance law between
two physical quantities of like dimension, and should not be read as one. This
caution is the reason the crossing is called a \emph{landmark} rather than an
equilibrium or a threshold in this section \cite{CubicThreshold}.
\end{remark}

% =====================================================================
\section{Domain Physics and Physical Embeddings}
\label{sec:embedding}
% =====================================================================

\begin{definition}[Domain embedding]
\label{def:embedding}
Let $X$ be a physical variable of a domain, and let the domain identify two
sectors whose normalized fractions are $a$ and $b$. A \emph{domain embedding}
is a map
\begin{equation}
\frac{a}{b}=F(X),
\qquad
a=\frac{F(X)}{1+F(X)},
\qquad
b=\frac{1}{1+F(X)} ,
\label{eq:embed}
\end{equation}
so that a trajectory in $X$ becomes a trajectory on the partition manifold.
\end{definition}

Two things must be supplied from outside the geometry before
\eqref{eq:embed} means anything.

First, the \emph{partition itself} must be genuine. The two sectors must
exhaust a conserved or otherwise fixed total for a reason internal to the
domain, not because they were divided by their sum for convenience. Any two
positive quantities can be normalized; a partition is physically meaningful
only when the normalization corresponds to a conservation law, a closure
relation, an exact decomposition, or a definitional completeness.

Second, the \emph{functional form of $F$} must be independently motivated. This
is the load-bearing requirement of the whole architecture, and \S\ref{sec:marker}
shows why: because the geometric landmark is fixed in partition space, the
physical location of that landmark is determined entirely by $F$. An
unconstrained $F$ can place a threshold wherever one wishes.

\begin{observation}[Division of labor]
\label{obs:labor}
Constraint geometry fixes the set of available relational landmarks and their
locations in partition space. Domain physics fixes which sectors are being
partitioned, the embedding $F$, the trajectory, and the interpretation. A
physical diagnostic exists only where both are supplied, and only survives
scrutiny when the second is independently justified.
\end{observation}

% =====================================================================
\section{The Perturbative General-Relativity Compactness Example}
\label{sec:gr}
% =====================================================================

This section develops the motivating example. It is the case in which the
partition is most clearly not imposed by hand, because the two sectors come
from an exact rearrangement of a field equation.

\subsection{An exact two-sector decomposition}

Write $g_{\mu\nu}=\bar g_{\mu\nu}+h_{\mu\nu}$ with $\bar g$ a background and
$h$ the perturbation. The Einstein tensor separates exactly into its part
linear in $h$ and the remainder,
\begin{equation}
G_{\mu\nu}[g]
= \underbrace{G^{\lin}_{\mu\nu}[h]}_{\text{linear}}
+ \underbrace{G^{\nonlin}_{\mu\nu}[h]}_{\text{higher order}} ,
\label{eq:Gsplit}
\end{equation}
so that the field equation may be rewritten as
\begin{equation}
G^{\lin}_{\mu\nu}[h]
= 8\pi G\,\bigl(T_{\mu\nu}+t_{\mu\nu}\bigr),
\qquad
t_{\mu\nu}\equiv -\frac{1}{8\pi G}\,G^{\nonlin}_{\mu\nu}[h] .
\label{eq:rearrange}
\end{equation}
Equation~\eqref{eq:rearrange} is a rearrangement, not an approximation: the
matter source and the gravitational self-energy contribution together form the
effective source of the linear operator. The object $t_{\mu\nu}$ is
gauge-dependent and is related in harmonic gauge to the standard
gravitational stress pseudotensor \cite{LL,ChoquetBruhat}.

This is what makes the two-sector reading legitimate here: the domain, not the
geometry, says that there are exactly two channels and that they exhaust the
effective source.

\subsection{The schematic sector scaling}

For a quasi-spherical configuration with compactness
\begin{equation}
C=\frac{2GM}{rc^{2}} ,
\label{eq:compactness}
\end{equation}
the perturbation scales as $|h|\sim C$, and the leading higher-order terms of
the Einstein tensor, being quadratic in $h$ and its derivatives, motivate the
schematic estimate
\begin{equation}
\frac{\text{nonlinear sector}}{\text{linear sector}}
\;=\;
\frac{\bigl\|G^{\nonlin}\bigr\|}{\bigl\|G^{\lin}\bigr\|}
\;\sim\; C .
\label{eq:scaling}
\end{equation}
Statement \eqref{eq:scaling} belongs to standard perturbation
theory \cite{MTW,ChoquetBruhat} and is not derived here. It is
schematic in two respects that matter later: it is an
order-of-magnitude relation between norms rather than an equality, and both
norms are gauge- and slicing-dependent.

\begin{remark}[Structural versus orbital compactness]
\label{rem:structural}
Two distinct quantities share the algebraic form \eqref{eq:compactness} and are
routinely both called compactness, and the distinction matters for anything said
later about perturbative accuracy. In \eqref{eq:compactness} as used here, $r$
is a radius of the configuration itself, so $C$ is a \emph{structural} quantity:
a fixed property of a star or a static geometry. In a compact binary, the
post-Newtonian expansion is instead controlled by the orbital parameter
$v^{2}/c^{2}\sim GM/(rc^{2})$ with $r$ the orbital \emph{separation}, a
\emph{dynamical} quantity that sweeps upward through the inspiral. The two are
not the same object, and a landmark value of one is not a landmark value of the
other. Consequently, an observation that some numerical value is close both to
the compactness of massive neutron stars and to the regime where post-Newtonian
accuracy degrades is not one coincidence but two claims about two quantities
that happen to share a symbol. This note keeps the structural reading
throughout, and makes no claim about the orbital one.
\end{remark}

\subsection{Normalization and the partition}

Normalizing the two sectors to a partition, and writing the ratio as the
embedding $F(C)=C$,
\begin{equation}
\frac{a}{b}=C,
\qquad
a=\frac{C}{1+C},
\qquad
b=\frac{1}{1+C},
\qquad
a+b=1 .
\label{eq:grembed}
\end{equation}
This is a domain embedding in the sense of Definition~\ref{def:embedding}. Its
elementary properties are: $a<b$ for $C<1$, so the higher-order sector is the
minority channel in the sub-horizon range; $a=b$ at $C=1$; and $a\to 0$ as
$C\to 0$, recovering a single-channel linear description in the weak field.

Independently of \eqref{eq:grembed}, the constraint geometry of
\S\ref{sec:geometry} supplies
\begin{equation}
h=\sqrt{ab}=\frac{\sqrt C}{1+C} .
\label{eq:grh}
\end{equation}

\subsection{The intersection}

Imposing the intrinsic crossing condition $h=a/b$ of
Proposition~\ref{prop:crossing} on the embedded partition gives, in partition
coordinates, $a=b^{3}$, and in the physical coordinate
\begin{equation}
\frac{\sqrt C}{1+C}=C
\quad\Longrightarrow\quad
C(1+C)^{2}=1
\quad\Longleftrightarrow\quad
C^{3}+2C^{2}+C-1=0 ,
\label{eq:cubicC}
\end{equation}
whose unique positive root is
\begin{equation}
\Ccrit=\psi-1=0.465571231876768\ldots=h_{*}=r_{*} .
\label{eq:Ccrit}
\end{equation}

\begin{proof}[Derivation of \eqref{eq:cubicC}]
From \eqref{eq:grh}, $h=C$ reads $\sqrt C=C(1+C)$; squaring gives
$C=C^{2}(1+C)^{2}$ and dividing by $C>0$ gives $C(1+C)^{2}=1$. Equivalently,
substituting \eqref{eq:grembed} into $b^{3}=a$ gives $(1+C)^{-3}=C(1+C)^{-1}$,
i.e.\ the same polynomial after multiplication by $(1+C)^{3}$.
\end{proof}

\begin{remark}[The two routes are one route]
\label{rem:notindependent}
The coincidence $\Ccrit=h_{*}$ in \eqref{eq:Ccrit} is algebraic, not
evidential. Both computations impose the same condition $a=b^{3}$ on the same
partition and differ only in which coordinate is solved for; the compactness
polynomial $C^{3}+2C^{2}+C-1$ and the altitude polynomial
$h^{3}+2h^{2}+h-1$ are the same polynomial. This equality is therefore not an
independent cross-check, and it is recorded here so that it is not mistaken
for one \cite{CubicThreshold}.
\end{remark}

\subsection{The division of labor in this example}

\begin{center}
\begin{tabular}{@{}ll@{}}
\toprule
Ingredient & Supplied by \\
\midrule
Two sectors exhausting the effective source & the exact rearrangement
\eqref{eq:rearrange} \\
The sector ratio $a/b=C$ & general relativity, via \eqref{eq:scaling} \\
The altitude $h=\sqrt{ab}$ & constraint geometry \\
The intrinsic condition $h=a/b$, i.e.\ $a=b^{3}$ & constraint geometry \\
The cubic $C(1+C)^{2}=1$ and its root & the intersection of the two \\
Any physical meaning of that root & neither, so far \\
\bottomrule
\end{tabular}
\end{center}

\begin{interpretation}
General relativity supplies the physical sector ratio. Constraint geometry
supplies the intrinsic consistency condition. Their intersection supplies a
cubic compactness diagnostic --- a candidate diagnostic, with a definite
number attached, and nothing more until the embedding is independently
justified and the diagnostic tested.
\end{interpretation}

\begin{remark}[Explicit disclaimers for this section]
Constraint geometry does not derive the Einstein field equations; the
uniqueness of $G_{\mu\nu}+\Lambda g_{\mu\nu}$ is a separate matter with a
separate proof \cite{Lovelock71,Wald}. The cubic \eqref{eq:cubicC} is not a
universal law of gravity. And $\Ccrit$ is not a point at which nonlinear
gravity begins: the higher-order sector is nonzero for every $C>0$ and grows
smoothly, so no onset is asserted. What \eqref{eq:cubicC} locates is where a
particular relational condition on a particular normalized decomposition is
met.

As to accuracy: the post-Newtonian expansion is reliable during the weak-field,
slow-motion inspiral and progressively loses accuracy as a binary approaches the
strong-field merger regime, where resummation methods or numerical-relativity
calculations become necessary \cite{Blanchet2024}. That statement is
qualitative, concerns the orbital parameter rather than the structural
compactness of \eqref{eq:compactness} (Remark~\ref{rem:structural}), and
attaches no numerical onset to any value. Nothing here asserts that accuracy
degrades near $\Ccrit$, and no such claim should be read into
\eqref{eq:cubicC}: establishing one would require a dedicated study specifying
the gauge, the waveform family, the post-Newtonian order, and the error measure
against which degradation is defined. Note also that ``loses accuracy'' is the
defensible claim and is weaker than ``diverges'': the post-Newtonian series is
an asymptotic expansion whose convergence properties are not established, so
what degrades is truncation accuracy at fixed order.
\end{remark}

% =====================================================================
\section{Intrinsic Marker Versus Physical Realization}
\label{sec:marker}
% =====================================================================

This section states the conceptual distinction that the rest of the note
depends on.

The intrinsic marker is fixed in partition space:
\begin{equation}
a=b^{3},\qquad a+b=1,
\qquad\text{hence}\qquad
r_{*}=\frac{a_{*}}{b_{*}}=\psi^{-2}=0.465571\ldots
\label{eq:rstar}
\end{equation}
Its physical realization, however, is the solution of
\begin{equation}
F(X_{*})=r_{*} ,
\label{eq:realize}
\end{equation}
which depends on the embedding. The geometry does not move. The physical
coordinate $X_{*}$ moves whenever $F$ changes.

Three distinct objects must therefore be kept apart:

\begin{enumerate}[label=\textbf{(\arabic*)},leftmargin=*]
\item \textbf{The intrinsic geometric marker.} The locus $a=b^{3}$ on the
partition, with its fixed values $(a_{*},h_{*},b_{*})=(\psi^{-3},\psi^{-2},\psi^{-1})$.
Exact, domain-free, and carrying no physical content
(Remark~\ref{rem:landmarkstatus}).

\item \textbf{The domain-specific embedding.} The map $a/b=F(X)$, which
converts the marker into a value $X_{*}$ of a physical variable via
\eqref{eq:realize}. This is where all the physics of the identification lives,
and where all the freedom to be wrong lives.

\item \textbf{The physical interpretation.} A claim that $X_{*}$ names
something --- a transition, a bound, a degradation of an approximation, an
observable clustering. This requires evidence beyond (1) and (2).
\end{enumerate}

\begin{observation}[The threshold is as firm as the embedding]
\label{obs:firmness}
Suppose the sector ratio is more accurately $a/b=\alpha C$ for some
dimensionless $\alpha$ not fixed by \eqref{eq:scaling}, which is an
order-of-magnitude statement between gauge-dependent norms. Then
\eqref{eq:realize} gives $C_{*}=r_{*}/\alpha$, and the physical threshold
scales inversely with a quantity the derivation never pinned down. More
generally, for $F(C)=\alpha C^{n}$ one obtains
$C_{*}=(r_{*}/\alpha)^{1/n}$. The intrinsic marker is exact; the number
$0.4656$ is exact \emph{as a partition ratio}; the compactness value inherits
the full indeterminacy of the normalization of \eqref{eq:scaling}.
\end{observation}

\begin{safeguard}[The manufacturing objection]
\label{safe:manufacture}
Given any target value $X_{0}$ and the fixed $r_{*}$, the embedding
$F(X)=r_{*}X/X_{0}$ places the landmark exactly at $X_{0}$. Therefore no
agreement between a partition landmark and an observed value is evidence for
anything unless $F$ was fixed, in full, by domain physics before the
comparison was made. This is the single most important safeguard in the
programme, and it is the reason the perturbative-GR case is the primary example:
there, $F(C)=C$ follows from an operator estimate that exists independently of
any partition consideration --- up to precisely the normalization ambiguity of
Observation~\ref{obs:firmness}, which must be reported rather than suppressed.
\end{safeguard}

% =====================================================================
\section{Relation to the Two-Face and Effective-Potential Frameworks}
\label{sec:priorframeworks}
% =====================================================================

The architecture \eqref{eq:architecture} is a refinement of an earlier
intuition, not a replacement for it. The progression can be described in three
stages.

\paragraph{Stage 1: the Two-Face Problem \cite{TwoFace}.}
Binary partition geometry was introduced as a translation framework. The same
constrained state admits a bounded description in $\chi$, a hyperbolic
description in $\xi=\operatorname{arctanh}\chi$, and a radial or response
description in $\Rresp=1/(1-\chi^{2})$; apparent conflicts between
descriptions of a system can be conflicts between faces rather than between
physical claims. In the present language, Stage 1 established that the
constraint layer possesses a well-defined coordinate vocabulary, and that
divergences in one face may be boundary effects rather than physical
singularities.

\paragraph{Stage 2: effective potentials on the partition manifold \cite{EffPot}.}
The Fisher metric of the binary partition,
\begin{equation}
ds_{F}^{2}=\frac{da^{2}}{a}+\frac{db^{2}}{b}
=\frac{d\chi^{2}}{1-\chi^{2}}
=\Rresp\,d\chi^{2},
\label{eq:fisher}
\end{equation}
was proposed as a universal kinetic skeleton, with a domain potential
$V_{\mathrm{domain}}(\chi)$ carrying the contextual dynamics. The same split
appears in the first-integral extraction as
\begin{equation}
V^{\lin}_{\mathrm{eff}}(\chi)
= -\tfrac12\,\frac{(\partial_{x}\chi)^{2}}{1-\chi^{2}}
= -\tfrac12\,
\frac{\text{contextual numerator}}{\text{universal denominator}} ,
\label{eq:numden}
\end{equation}
which is the dynamical instance of the architecture:
\begin{equation}
\begin{gathered}
\text{universal Fisher denominator}\;+\;\text{contextual numerator}\\
\longleftrightarrow\\
\text{universal constraint structure}\;+\;\text{domain-specific dynamics}.
\end{gathered}
\end{equation}
Stage 2 also established the discipline that the present note inherits: applied
to a single known trajectory, the extraction \eqref{eq:numden} is tautological,
because the free function $V$ has as many degrees of freedom as the data.
Single-trajectory extraction is therefore a Tier-1 normal form and acquires
falsifiable status only after an over-determination test --- a shared-$V$ family
test with a reported residual, an independently determined $V$ that then
predicts the trajectory, cross-domain recurrence under a pre-declared
normalization, or genuine compression \cite{EffPot}.

\paragraph{Stage 3: perturbative GR and the intrinsic cubic
\cite{CubicThreshold,IntrinsicCubic,ReverseThales}.}
The constraint layer was shown to supply not only coordinates and a kinetic
skeleton but also intrinsic relational \emph{landmarks} --- loci fixed by the
constraint alone --- with general relativity supplying an embedding under which
one such landmark acquires a compactness reading.

\begin{interpretation}[Refinement, not abandonment]
The original intuition was that a single constrained geometry underlies
structurally similar relations across domains. That intuition survives, in a
sharpened form: what recurs is the \emph{state space}, its coordinate
vocabulary, and its landmark structure --- not the dynamics, not the
mechanisms, and not the physical meanings. The mature formulation is:

\begin{quote}
\emph{The partition framework does not generate domain physics. It provides a
common constrained state space, a coordinate language, and a set of intrinsic
relational landmarks against which domain-specific dynamics can be organized
and tested.}
\end{quote}
\end{interpretation}

\begin{remark}[A coordinate caution carried forward]
The flat coordinate of the Fisher metric \eqref{eq:fisher} is $u=\arcsin\chi$,
not the rapidity $\xi$, which is flat for the distinct hyperbolic metric
$d\chi^{2}/(1-\chi^{2})^{2}$. Writing a Fisher-denominator kinetic term in $\xi$
therefore leaves a residual $\operatorname{sech}^{2}\xi$ weight that is a
signature of the metric--coordinate pairing and not of domain physics
\cite{EffPot}. Comparisons across domains must declare the pairing.
\end{remark}

% =====================================================================
\section{Cross-Domain Illustrations}
\label{sec:crossdomain}
% =====================================================================

These illustrations are deliberately brief. Their purpose is to show that the
same intrinsic marker occurs in each partition while its physical status
differs completely --- which is exactly the point of the architecture. The
perturbative-GR case of \S\ref{sec:gr} remains the central example.

\subsection{Flat $\Lambda$CDM: two landmarks, failing in opposite ways}
\label{sec:lcdm}

In a flat cosmology the density parameters exhaust unity, so the partition is
genuine rather than imposed: the closure is a constraint of the Friedmann
system, not a normalization of convenience. In the matter-plus-$\Lambda$ regime
where radiation is negligible, take
\begin{equation}
a=\Omega_{m},\qquad b=\Omega_{\Lambda},\qquad a+b=1 ,
\label{eq:lcdmassign}
\end{equation}
so that the partition geometry organizes the coupling
$\mu=ab=\Omega_{m}\Omega_{\Lambda}$, the asymmetry
$\chi=\Omega_{m}-\Omega_{\Lambda}$, and the response $R=1/(1-\chi^{2})$.

\begin{remark}[Which sector is $a$, and why it matters]
\label{rem:ordering}
The condition $a=b^{3}$ is not symmetric under exchange of the sectors, whereas
$h=\sqrt{ab}$, $\chi^{2}$ and $R$ are. Identifying a landmark by its response
value therefore specifies an unordered \emph{pair} of states: $R=1.15337\ldots$
is attained both where $\Omega_{m}=\Omega_{\Lambda}^{3}$ and where
$\Omega_{\Lambda}=\Omega_{m}^{3}$, and these are different epochs. Only the
ordered condition selects a branch, and only domain physics can order the
sectors. The assignment \eqref{eq:lcdmassign} is the one under which the
occupied state is the marker: matter is at present the minority sector, with
$\Omega_{m}\approx 0.317\approx a_{*}$, which is also the orientation of the
general-relativity embedding of \S\ref{sec:gr}, where the higher-order sector
is the minority channel. It additionally aligns the sign convention, since
$\chi_{*}<0$ in Proposition~\ref{prop:crossing} and
$\chi_{\mathrm{today}}=\Omega_{m}-\Omega_{\Lambda}<0$. Reversing the assignment
moves the marker to a substantially earlier epoch and flips the sign of $\chi$;
nothing in the constraint geometry chooses between the two orderings.
\end{remark}

The embedding is supplied by the conservation laws rather than chosen. With
scale factor $s$, matter dilution $\rho_{m}\propto s^{-3}$ and
$\rho_{\Lambda}$ constant give
\begin{equation}
\frac{a}{b}=F(s)=\frac{\Omega_{m,0}}{\Omega_{\Lambda,0}}\,s^{-3} ,
\label{eq:lcdmembed}
\end{equation}
and, more generally, the rapidity
$\xi_{\mathrm{DE}}=\tfrac12\ln(\rho_{\mathrm{DE}}/\rho_{m})$ obeys
\begin{equation}
\frac{d\xi_{\mathrm{DE}}}{d\ln s}
=-\tfrac{3}{2}\,w_{\mathrm{DE}}(s) ,
\label{eq:rapidityw}
\end{equation}
which follows from the two separate conservation laws alone --- $w_{m}=0$ and
$\rho_{\mathrm{DE}}'/\rho_{\mathrm{DE}}=-3(1+w_{\mathrm{DE}})$ --- and requires
neither flatness nor two-component closure \cite{PartitionFlow}. This is an
embedding of the stronger kind: it is the independent-determination route of
\cite{EffPot}, with $\Lambda$CDM the case $w_{\mathrm{DE}}=-1$, for which the
flow is a uniform translation in rapidity at rate $3/2$.

Two landmarks are now available on this trajectory, and they fail in opposite
ways.

\paragraph{(a) A landmark with a mechanism and no occupancy.}
The domain flow supplies its own landmark. Writing $J=\mu|\chi|$, flat
$\Lambda$CDM satisfies exactly
\begin{equation}
\frac{d\Omega_{m}}{d\ln s}=-3\mu,
\qquad
\frac{d^{2}\Omega_{m}}{(d\ln s)^{2}}=9\mu\chi=9J\operatorname{sgn}\chi ,
\label{eq:flowid}
\end{equation}
so the coupling is the velocity of the matter-fraction flow and the
coupling--asymmetry product is its acceleration; the acceleration extremum lies
at $|\chi|=1/\sqrt3$, i.e.\ at $R_{J}=3/2$, equivalently at the state where
$d^{3}\Omega_{m}/(d\ln s)^{3}$ vanishes \cite{PartitionFlow}. This landmark is
\emph{derived from the domain mechanism} and therefore means something about the
flow. The present epoch is not there: $R_{J}$ is attained at $z\approx 1.0$ and
$z\approx -0.17$, some twenty standard deviations from today in the
cosmological parameters \cite{PartitionFlow,Planck18}.

\paragraph{(b) A landmark with occupancy and no mechanism.}
The intrinsic marker $a=b^{3}$ is realized at $F(s_{*})=r_{*}$, i.e.\ at
\begin{equation}
s_{*}=\Bigl(\frac{\Omega_{m,0}}{r_{*}\,\Omega_{\Lambda,0}}\Bigr)^{1/3} ,
\end{equation}
which for concordance parameters falls at $z\approx 0.002$ --- effectively now,
with $R_{*}=1.15337$ against a present-day $R\approx 1.156$
\cite{PartitionFlow,Planck18}. Here occupancy is excellent and the mechanism is
absent; and in this domain the absence can be \emph{proved} rather than merely
conceded. Using $d\chi/d\ln s=\tfrac32(1-\chi^{2})$, the flow admits an exact
Fisher-gradient representation
\begin{equation}
\frac{d\chi}{d\ln s}=-(1-\chi^{2})\,U'(\chi),
\qquad
U(\chi)=-\tfrac{3}{2}\chi ,
\label{eq:linearU}
\end{equation}
in which the potential is linear, so $U'\equiv-\tfrac32\neq 0$ everywhere. The
flow therefore has no interior critical point at all: its only fixed points are
the simplex edges $\chi=\pm1$, and at the cubic state
$d\chi/d\ln s=\tfrac32(1-\chi_{*}^{2})\approx 1.30\neq 0$ \cite{PartitionFlow}.
The universe passes through the cubic state at ordinary speed. A mechanism that
selected it would require a domain potential with an interior minimum there,
which is not $\Lambda$CDM.

\begin{interpretation}[The sharpest form of the thesis]
\label{int:sharpest}
Neither cosmological landmark is a diagnostic, and they miss for opposite
reasons. Landmark (a) is derived from the domain and unoccupied; landmark (b)
is intrinsic to the constraint geometry and occupied but provably mechanism-free.
Meaning and occupancy are supplied by different layers, and neither implies the
other. A framework that reported only (b) would look like a success; a
framework that reported only (a) would look like a mechanism. Reporting both is
what makes the architecture of \eqref{eq:architecture} visible.
\end{interpretation}

\subsection{The effective-potential framework}

Under a pullback $\chi=\chi_{\mathrm{domain}}(x)$ the same split reappears
dynamically as \eqref{eq:numden}: the denominator $1-\chi^{2}$ is contributed
by the chart, and the numerator $(\partial_{x}\chi)^{2}$ by the domain approach
law. Shared boundary divergence as $|\chi|\to 1$ is therefore a shared response
geometry and not shared physics; two systems can share the boundary class while
occupying different numerator classes \cite{EffPot}. Single-trajectory
extraction remains a Tier-1 normal form unless it passes an over-determination
test, as recalled in \S\ref{sec:priorframeworks}.

\subsection{Newtonian two-body geometry}

Fixed total mass gives a genuine partition. With $c=m_{1}+m_{2}=M$ and
segments $m_{1},m_{2}$, the altitude theorem gives $h^{2}=m_{1}m_{2}$ exactly,
so the gravitational force and potential energy are altitude-squared
quantities, $F_{\mathrm{grav}}=Gh^{2}/r^{2}$ and $U=-Gh^{2}/r$, while the
reduced mass is $\mu=h^{2}/c$ and the symmetric mass ratio is $\nu=h^{2}/c^{2}$
\cite{Newtonian2Body}. In normalized form $a=m_{1}/M$, $b=m_{2}/M$ the
embedding is the mass ratio itself, $F(q)=q$ with $q=m_{1}/m_{2}$, so the
intrinsic marker is realized at the single mass ratio $q_{*}=r_{*}$, i.e.\
$m_{2}/m_{1}=1/r_{*}=\psi^{2}=2.1479\ldots$

This is the most instructive of the three illustrations precisely because
nothing happens there. Newtonian mechanics supplies the physical laws and the
observables; it assigns no distinguished role to that mass ratio. The intrinsic
marker exists in the partition, is exactly located, and is physically empty in
this domain. Meanwhile the same geometry does useful organizational work
elsewhere in the same system: it exhibits the coupling--timing asymmetry
whereby force is altitude-governed and Kepler timing is hypotenuse-governed,
which is an exact re-expression of standard relations and not a new result
\cite{Newtonian2Body}.

\begin{observation}[Landmark availability versus landmark realization]
Across the illustrations the marker $a=b^{3}$ is available in every case and
means something different in each: a candidate diagnostic under a schematic
operator embedding (\S\ref{sec:gr}); an occupied but provably mechanism-free
transit state under an independently derived evolution law
(\S\ref{sec:lcdm}(b)); and nothing at all under a trivially exact embedding
(two-body masses). Alongside it, $\Lambda$CDM carries a second landmark that is
mechanism-bearing and unoccupied (\S\ref{sec:lcdm}(a)). Availability is
geometric; realization is physical; and the two can be present separately.
\end{observation}

% =====================================================================
\section{A Protocol, and its Safeguards}
\label{sec:protocol}
% =====================================================================

The following is a conservative working protocol. Steps 1--4 are the
constructive part; step 5 is what makes the result scientific rather than
decorative.

\begin{enumerate}[label=\textbf{Step \arabic*.},leftmargin=*,itemsep=4pt]
\item \textbf{Identify a genuine constrained decomposition.} Locate two
sectors that exhaust a conserved total, an exact operator decomposition, or a
definitionally complete budget, for reasons internal to the domain. Reject
partitions created by dividing two unrelated quantities by their sum.

\item \textbf{Normalize to $a+b=1$.} Record the normalization explicitly,
including any gauge, slicing, or norm choice on which it depends.

\item \textbf{Derive the domain map.} Obtain $a/b=F(X)$ from domain physics ---
a conservation law, an evolution equation, an operator estimate, an equation of
state --- and state which parts of $F$ are derived and which are conventional.
In particular, report any undetermined overall normalization, since by
Observation~\ref{obs:firmness} it propagates directly into the location of any
landmark.

\item \textbf{Evaluate the landmarks.} Solve $F(X_{*})=r_{*}$ for the intrinsic
marker of interest, and ask whether $X_{*}$ is a physically meaningful
coordinate: does the domain independently distinguish it? Is it inside the
physically admissible range? Does it correspond to any observable change?

\item \textbf{Require independent justification or falsifiable testing of the
embedding.} An arbitrary choice of $F$ can manufacture any desired threshold
(Safeguard~\ref{safe:manufacture}). The embedding must therefore either follow
from independent domain physics, or the resulting diagnostic must make a
prediction that could fail: a held-out prediction, a collapse of a family of
trajectories onto one structure with a reported residual, a cross-domain
recurrence under a pre-declared normalization, or a compression of a complex
description into a shorter one \cite{EffPot}.
\end{enumerate}

\subsection{Additional safeguards}

\begin{safeguard}[Universality trap]
Many systems normalize to $a+b=1$. Shared arithmetic is not shared physics, and
the recurrence of the partition across domains is expected rather than
informative.
\end{safeguard}

\begin{safeguard}[Chart contributions]
Divergences and boundary behaviour contributed by the coordinates ---
$\Rresp\to\infty$ as $|\chi|\to1$, the Fisher denominator, the response face
generally --- must not be reported as domain findings.
\end{safeguard}

\begin{safeguard}[Algebraic identities are not cross-checks]
When two computations impose the same condition on the same partition and
differ only in which coordinate is solved for, their agreement is a tautology
(Remark~\ref{rem:notindependent}). Cross-checks must involve independent
inputs.
\end{safeguard}

\begin{safeguard}[Admissible range]
A partition coordinate normalized to $[0,1]$ need not have a physically
admissible range that large. Any landmark must be located relative to the
domain's own bounds, and a landmark outside them is not a diagnostic.
\end{safeguard}

\begin{safeguard}[Landmark crowding]
\label{safe:crowding}
Before reporting a landmark as occupied, check whether its neighbourhood
contains generic constants within the measurement error of the domain. Near the
cubic value the interval of width comparable to current cosmological precision
contains at least three unrelated numbers: the cubic root
$b_{*}=0.682328$; the Gaussian one-sigma mass
$\operatorname{erf}(1/\sqrt2)=0.682689$, a generic aggregation landmark with no
domain-specific content \cite{NormNulls}; and the present-day
$\Omega_{\Lambda,0}\approx 0.6835$ with an uncertainty of order $0.007$
\cite{Planck18,PartitionFlow}. The separations are an order of magnitude below
the uncertainty, so proximity cannot discriminate among them and agreement with
any one of them is not structural. The safeguard is general: a landmark is only
as informative as its neighbourhood is sparse.
\end{safeguard}

\begin{safeguard}[Information-content check]
\label{safe:infocontent}
Before testing whether a partition coordinate organizes an observable, check
whether it can. If every partition coordinate is a deterministic invertible
function of one domain quantity, then all of them carry exactly the information
already in that quantity and none adds a predictive degree of freedom. In flat
$\Lambda$CDM, $\mu$, $\chi$, $R$, $J$ and $\xi$ are all functions of
$\Omega_{m}(z)$ alone, so any apparent value added by a higher-order coordinate
over a lower-order one signals leakage --- residual dependence re-entering
through a nonlinear reparameterization --- rather than new content. Genuine
content requires leaving the regime in which the reduction holds; for the
cosmological case that means $w_{\mathrm{DE}}\neq-1$, where the rapidity
curvature $d^{2}\xi/(d\ln s)^{2}=-\tfrac32 w_{\mathrm{DE}}'$ is nonzero
\cite{PartitionFlow}. This is a screen to be applied before a test is run, not
an explanation offered after one fails.
\end{safeguard}

\begin{safeguard}[Order and level]
Conditions relating an amplitude-level quantity such as $h$ to a ratio-level
quantity such as $a/b$ (Remark~\ref{rem:levels}) must be labeled as
partition-intrinsic relational conditions rather than physical balance laws.
\end{safeguard}

% =====================================================================
\section{Epistemic Tiers}
\label{sec:tiers}
% =====================================================================

\paragraph{Tier 1 --- exact mathematics.}
\begin{itemize}[leftmargin=*,itemsep=1pt]
\item the constraint $a+b=1$ with $a,b>0$;
\item the altitude $h=\sqrt{ab}$ and the identity $4h^{2}+\chi^{2}=1$;
\item the response and rapidity coordinates and all coordinate transformations
      among the coordinates $(a,b,h,\chi,\Rresp,\xi)$;
\item the intrinsic crossing $h=a/b$, its equivalence to $a=b^{3}$, its
      uniqueness, and the sign law of Lemma~\ref{lem:sign};
\item the resulting cubic relations $b^{3}+b-1=0$ and $h^{3}+2h^{2}+h-1=0$,
      the supergolden coordinates
      $(a_{*},h_{*},b_{*})=(\psi^{-3},\psi^{-2},\psi^{-1})$, and the algebraic
      identity of the compactness and altitude polynomials;
\item the algebraic consequence \eqref{eq:cubicC} of imposing $a=b^{3}$ under
      the embedding $a/b=C$;
\item the Fisher metric \eqref{eq:fisher} and the numerator/denominator split
      \eqref{eq:numden}.
\end{itemize}

\paragraph{Tier 2 --- domain diagnostic.}
\begin{itemize}[leftmargin=*,itemsep=1pt]
\item an independently motivated identification of the two physical sectors,
      including the reason they exhaust a total;
\item a derived map $a/b=F(X)$, with its conventional content declared;
\item the physical location $X_{*}$ of a geometric landmark, reported together
      with its sensitivity to the normalization of $F$;
\item a falsifiable test: a held-out prediction, a family collapse with a
      reported residual, an independently determined structure that then
      predicts a trajectory, or cross-domain recurrence under pre-declared
      normalization.
\end{itemize}
Nothing in the perturbative-GR example reaches Tier 2 in this note: the
embedding \eqref{eq:grembed} is schematic (Observation~\ref{obs:firmness}) and
no test is performed here.

\paragraph{Tier 3 --- speculative interpretation.}
\begin{itemize}[leftmargin=*,itemsep=1pt]
\item claims of a universal mechanism behind the recurrence of the partition;
\item claims that the intrinsic cubic is fundamental, or that its root is a
      constant of nature;
\item claims that distinct physical domains share dynamics because they share
      partition geometry;
\item claims that a landmark marks a sharp physical onset.
\end{itemize}

% =====================================================================
\section{Discussion}
\label{sec:discussion}
% =====================================================================

Three points deserve emphasis.

First, the architecture explains why the programme's strongest and weakest
results look superficially alike. Both consist of a partition, an embedding,
and a landmark. They differ in exactly one place: whether the embedding was
derived or chosen. The perturbative-GR case is the strongest available because
the decomposition \eqref{eq:rearrange} is exact and the ratio
\eqref{eq:scaling} pre-exists the partition; it remains short of a diagnostic
because \eqref{eq:scaling} is schematic. The two-body case is the weakest not
because its embedding is doubtful --- it is exact --- but because the domain
assigns the landmark no role.

Second, the architecture makes the failure modes explicit and separable. If a
partition is not genuine, step~1 fails and nothing downstream can be repaired.
If the embedding is chosen rather than derived, step~5 fails and any numerical
agreement is manufactured. If the landmark is real and the embedding is derived
but the domain distinguishes nothing there, the result is a transit marker: a
correct, exactly located, physically empty statement. These are three different
negative results, and distinguishing them is more useful than aggregating them
into a single verdict.

Third, the cases now on the table separate cleanly along these axes.

\begin{center}
\small
\begin{tabular}{@{}p{3.0cm}p{3.4cm}p{3.2cm}p{3.0cm}@{}}
\toprule
Case & Embedding & Domain role for landmark & Occupancy \\
\midrule
Perturbative GR, $\Ccrit$ (\S\ref{sec:gr})
& schematic: $a/b\sim C$ up to an unpinned normalization
  (Obs.~\ref{obs:firmness})
& not established
& untested here \\
$\Lambda$CDM, $R_{J}$ (\S\ref{sec:lcdm}(a))
& derived from the conservation laws
& yes: extremum of the flow acceleration
& no ($z\approx1.0$, $-0.17$) \\
$\Lambda$CDM, cubic (\S\ref{sec:lcdm}(b))
& derived from the conservation laws
& provably none: linear $U(\chi)$
& yes ($z\approx0.002$) \\
Two-body, cubic (\S\ref{sec:crossdomain})
& exact: $F(q)=q$
& none
& not applicable \\
\bottomrule
\end{tabular}
\end{center}

Fourth, a comparison drawn in the earlier cosmological work should be read the
other way round. That work contrasted the linear cosmological potential
\eqref{eq:linearU} with a compactness flow in which the cubic is an attracting
state, which invites the reading that $\Lambda$CDM is the weaker case. Under the
present architecture it is the stronger one. The cosmological embedding is
derived --- \eqref{eq:rapidityw} follows from the conservation laws alone --- and
it returns a clean, provable negative. The compactness flow's potential, by
contrast, is a normal form obtained by quadrature and carries no information
beyond the sign field of the flow; its critical point is inherited from the
partition geometry rather than derived from general relativity, and the current
form of that result asserts robustness across a class of positive prefactors
rather than selection \cite{AltitudeRelax}. The two cases therefore divide the
labor differently: $\Lambda$CDM has the better embedding and an empty landmark,
while perturbative GR has the more suggestive landmark and a schematic embedding
with an unpinned normalization. Neither is a diagnostic, and the useful
statement is the pair rather than either one alone.

Fifth, negative results of the third kind are informative in the same way that
a failed family test is informative in the dynamical setting \cite{EffPot}. A
domain in which every intrinsic landmark is a transit marker is a domain whose
physics is not organized by its partition structure; that is a fact about the
domain, stated in a common language, and it is worth recording.

What the framework offers, then, is not derivation but organization with
declared seams: a fixed constrained state space, a coordinate vocabulary that
translates between bounded, hyperbolic and response descriptions, a
short list of exactly located intrinsic landmarks, and a protocol that makes
the physical content --- and the burden of proof --- reside entirely in the
embedding.

% =====================================================================
\section{Conclusion}
\label{sec:conclusion}
% =====================================================================

A conserved two-sector system may admit a universal partition geometry, but
that geometry does not determine the system's dynamics. It supplies invariant
coordinates, coupling measures, asymmetry measures, and intrinsic relational
landmarks. Domain physics supplies the sectors, the embedding, the evolution,
and the empirical meaning. A physical diagnostic emerges only where these two
layers intersect under an independently justified map.

This is also why the original construction required both ingredients at once,
and why neither could have been dispensed with. The Einstein equations grounded
the partition in a physically meaningful linear--nonlinear decomposition, while
the Thales semicircle supplied the intrinsic relational condition that selected
the cubic landmark. Without the field equations there would have been no reason
to regard the two sectors as a genuine partition of anything; without the
semicircle there would have been no condition to impose on it, and the operator
scaling alone yields only the trivial bound $C<1$. The cubic diagnostic is the
intersection, and it is no firmer than the weaker of the two ingredients --- at
present, the schematic normalization of the sector ratio.

% =====================================================================
\section{Relation to Prior Work}
\label{sec:priorwork}
% =====================================================================

This note synthesizes and clarifies results already developed in earlier work;
it introduces no new physical law, no new theorem, and no new empirical claim.
Specifically:

\begin{itemize}[leftmargin=*]
\item The coordinate vocabulary and the multi-face translation between bounded,
hyperbolic and response descriptions were developed in
\cite{TwoFace}.

\item The Fisher kinetic skeleton, the universal-denominator /
contextual-numerator split, the normal-form status of single-trajectory
extraction, and the over-determination routes cited in
\S\ref{sec:priorframeworks} and \S\ref{sec:protocol} were developed in
\cite{EffPot}.

\item The altitude--ratio crossing, its equivalence to $a=b^{3}$, the
uniqueness of the interior solution, and the supergolden coordinates of
Proposition~\ref{prop:crossing} were established in
\cite{IntrinsicCubic}.

\item The compactness embedding $a/b=C$, the operator rearrangement
\eqref{eq:rearrange}, and the cubic \eqref{eq:cubicC} were developed in
\cite{CubicThreshold}; the transversality of the crossing and the
prefactor-independence of the associated relaxation class, together with the
recalibration of language from ``dynamically selected'' to ``attracting
equilibrium of a modelling class,'' also appear
there.

\item The similarity-composition reading of the structural content of the field
equations, and the explicit statement that the field equations are not derived
from the geometry, appear in
\cite{ReverseThales}.

\item The exact flat-$\Lambda$CDM flow identities \eqref{eq:flowid}, the
landmark $R_{J}=3/2$, the rapidity identity \eqref{eq:rapidityw}, the
Fisher-gradient representation \eqref{eq:linearU} with its linear potential, and
the numerical redshift placement of all cosmological landmarks were established
in \cite{PartitionFlow}. The present note adds only the reading of that paper's
two landmarks as opposite failure modes
(Interpretation~\ref{int:sharpest}), and the ordering correction of
Remark~\ref{rem:ordering}.

\item The numerical crowding of the region near $0.68$ was recorded in
\cite{NormNulls}; Safeguard~\ref{safe:crowding} generalizes it beyond
cosmology.

\item \textbf{Citation hygiene.} The June 2026 note \cite{DynSelSuperseded} is
superseded by \cite{AltitudeRelax}, which withdraws the $C(1-C)$ prefactor and
replaces ``dynamically selected'' with ``attracting equilibrium of the
imbalance-driven class.'' Forward citations should point to the revised note, so
that the withdrawn artifacts are not carried into later
work.

\item \textbf{Second hygiene item.} The Blanchet living review has been
superseded by its current edition \cite{Blanchet2024}; earlier notes in this
programme cite the 2014 edition, and forward citations should be repointed. More
importantly, one of those notes attributes to that review a numerical statement
--- that post-Newtonian accuracy degrades for compactness above roughly
$0.4$--$0.5$ --- which the review does not supply; the source supports the
qualitative loss of accuracy approaching merger and no threshold value. That
attribution should be weakened to the qualitative form or
removed.

\item The mass-partition organization of Newtonian two-body quantities, the
altitude theorem $h^{2}=m_{1}m_{2}$, and the coupling--timing asymmetry were
developed in \cite{Newtonian2Body}.
\end{itemize}

The contribution of the present note is organizational: it names the two layers,
states the division of labor \eqref{eq:architecture}, separates the intrinsic
marker from its embedding and its interpretation (\S\ref{sec:marker}), and
records the safeguard that an arbitrary embedding can manufacture any threshold
(Safeguard~\ref{safe:manufacture}). Where earlier work stated these points
locally, in the context of a particular calculation, they are here stated once,
generally, and in a form that applies before the next calculation is begun.

% =====================================================================
\section*{Acknowledgment}
% =====================================================================
\addcontentsline{toc}{section}{Acknowledgment}
The author acknowledges the use of artificial-intelligence tools for editorial
assistance, structural organization, algebraic checking, and consistency
review. All conceptual development, research direction, interpretation, and
final responsibility for the content belong to the author.

% =====================================================================
\begin{thebibliography}{99}
% ---------------------------------------------------------------------
% Author's own work.
% ---------------------------------------------------------------------
\bibitem{TwoFace}
E. De Jes\'us, \emph{The Two-Face Problem: A Practitioner's Guide to Bridging
Inconsistencies Across Physics}, Zenodo (2026).
\href{https://doi.org/10.5281/zenodo.20221967}{DOI:10.5281/zenodo.20221967}.

\bibitem{EffPot}
E. De Jes\'us, \emph{Two-Face Series: Effective Potentials on the Binary
Partition Manifold}, Zenodo (2026).
\href{https://doi.org/10.5281/zenodo.21430864}{DOI:10.5281/zenodo.21430864}.

\bibitem{IntrinsicCubic}
E. De Jes\'us, \emph{The Partition-Intrinsic Cubic: The Unique Altitude--Ratio
Crossing, Supergolden Coordinates, and Algebraic Evaluation on the Thales
Partition Manifold}, Version 3, Zenodo (2026).
\href{https://doi.org/10.5281/zenodo.21543445}{DOI:10.5281/zenodo.21543445}.

\bibitem{CubicThreshold}
E. De Jes\'us, \emph{The Cubic Compactness Threshold in Perturbative General
Relativity: A Conservative Reverse-Thales Diagnostic of Linear--Nonlinear
Operator Balance}, Zenodo (2026).
\href{https://doi.org/10.5281/zenodo.20564050}{DOI:10.5281/zenodo.20564050}.

\bibitem{ReverseThales}
E. De Jes\'us, \emph{Iterated Similarity Composition and the Structure of the
Einstein Field Equations: A Reverse Thales Construction}, Zenodo (2026).

\bibitem{PartitionFlow}
E. De Jes\'us, \emph{Matter--Dark-Energy Balance as a Partition Flow: Velocity,
Acceleration, and Constant Rapidity in Flat $\Lambda$CDM}, Zenodo (2026).

\bibitem{AltitudeRelax}
E. De Jes\'us, \emph{Altitude--Ratio Relaxation on the Thales Partition
Manifold: the Cubic Compactness Threshold as a Robust Interior Attractor for a
Class of Imbalance-Driven Flows}, revised technical note, Zenodo (2026).
Supersedes Ref.~\cite{DynSelSuperseded}.

\bibitem{DynSelSuperseded}
E. De Jes\'us, \emph{Dynamical Selection of the Cubic Compactness Threshold:
Altitude--Ratio Imbalance, Partition-Intrinsic Flow, and the Stable Attractor},
Zenodo (June 2026).
\href{https://doi.org/10.5281/zenodo.20603097}{DOI:10.5281/zenodo.20603097}.
Superseded by Ref.~\cite{AltitudeRelax}, in which the $C(1-C)$ prefactor and
the ``dynamically selected'' framing are withdrawn.

\bibitem{NormNulls}
E. De Jes\'us, \emph{Two Normalization Nulls and the Tests That Tell Them
Apart}, Zenodo (2026).
\href{https://doi.org/10.5281/zenodo.20774968}{DOI:10.5281/zenodo.20774968}.

\bibitem{Newtonian2Body}
E. De Jes\'us, \emph{Newtonian Two-Body Gravity on the Thales Semicircle: A
Geometric Diagnostic for Mass-Partition Structure}, Zenodo (2026).
\href{https://doi.org/10.5281/zenodo.18496501}{DOI:10.5281/zenodo.18496501}.

% ---------------------------------------------------------------------
% EXTERNAL LITERATURE
% External literature.
% ---------------------------------------------------------------------
\bibitem{Wald}
R. M. Wald, \emph{General Relativity}, University of Chicago Press (1984).

\bibitem{MTW}
C. W. Misner, K. S. Thorne and J. A. Wheeler, \emph{Gravitation},
W. H. Freeman (1973).

\bibitem{Lovelock71}
D. Lovelock, ``The Einstein tensor and its generalizations,''
\emph{J. Math. Phys.} \textbf{12}, 498 (1971).

\bibitem{Blanchet2024}
L. Blanchet, ``Post-Newtonian theory for gravitational waves,''
\emph{Living Reviews in Relativity} \textbf{27}, 4 (2024).
\href{https://doi.org/10.1007/s41114-024-00050-z}{DOI:10.1007/s41114-024-00050-z}.

\bibitem{LL}
L. D. Landau and E. M. Lifshitz, \emph{The Classical Theory of Fields}, 4th ed.,
Butterworth--Heinemann (1975).

\bibitem{ChoquetBruhat}
Y. Choquet-Bruhat, \emph{General Relativity and the Einstein Equations},
Oxford University Press (2009).

\bibitem{Amari}
S. Amari and H. Nagaoka, \emph{Methods of Information Geometry}, AMS (2000).

\bibitem{OEIS}
OEIS Foundation Inc., \emph{A000930: Narayana's cows sequence}, The On-Line
Encyclopedia of Integer Sequences, \url{https://oeis.org/A000930}.

\bibitem{Planck18}
Planck Collaboration (N.~Aghanim \emph{et al.}), ``Planck 2018 results.~VI.
Cosmological parameters,'' \emph{Astron. Astrophys.} \textbf{641}, A6 (2020).
\href{https://doi.org/10.1051/0004-6361/201833910}{DOI:10.1051/0004-6361/201833910}.

\end{thebibliography}

\end{document}
